\section{Primitive comparison with \texorpdfstring{\(\mathfrak g_{\Delta_5}\)}{gDelta5}: attack--heal}
\label{sec:primitive-comparison-attack-heal}

\providecommand{\Z}{\mathbb Z}
\providecommand{\C}{\mathbb C}
\providecommand{\La}{\Lambda}
\providecommand{\Prim}{\operatorname{Prim}}
\providecommand{\sdim}{\operatorname{sdim}}
\providecommand{\smult}{\operatorname{smult}}
\providecommand{\autcor}{\mathfrak g_{\Delta_5}}
\providecommand{\Hall}{\mathcal H}
\providecommand{\BM}{\operatorname{BM}}
\providecommand{\Fact}{\mathsf{Fact}}
\providecommand{\Den}{\mathsf{Den}}

\subsection{Claim attacked}

The attacked implication is
\[
\begin{gathered}
\hbox{compact }E_3/\hbox{Hall/factorization source for }K3\times E,\\
\sdim H^0\Prim_{\mathrm{prot}}(\mathcal F_X)_{(n,l,m)}=f(nm,l),\\
\operatorname{den}(\autcor)=64^{-1}\Delta_5(2Z)
\end{gathered}
\quad\Longrightarrow\quad
H^0\Prim_{\mathrm{prot}}(\mathcal F_X)\cong\autcor .
\]
This implication is not proved.  The manuscript proves the virtual
determinant and the Gritsenko--Nikulin/Borcherds denominator algebra.  It
does not construct the primitive Hall/factorization bracket, nor the map
identifying that bracket with the Borcherds bracket.

Local anchors read before this note:
\[
\begin{array}{ll}
\texttt{main.tex:250--296} & \hbox{primitive bracket and comparison marked as realization problem},\\
\texttt{main.tex:496--522} & \hbox{virtual determinant has no products or locality},\\
\texttt{main.tex:1566--1671} & \hbox{definition and proof of the denominator algebra identity},\\
\texttt{main.tex:1907--2003} & \hbox{compact Igusa realization datum and primitive comparison},\\
\texttt{main.tex:2059--2082} & \hbox{determinant does not determine Hall constants},\\
\texttt{main.tex:2117--2243} & \hbox{Borcherds simple-root parity, relations, denominator},\\
\texttt{main.tex:2283--2380} & \hbox{what the denominator and determinant forget},\\
\texttt{agent\_material/10\_bps\_bracket\_constants\_attack\_heal.tex} &
\hbox{low-degree constants and first Hall gap},\\
\texttt{agent\_material/11\_factorization\_koszul\_minimal\_data.tex} &
\hbox{five-object compact realization dossier},\\
\texttt{agent\_material/08\_compact\_hall\_factorization\_object.tex} &
\hbox{compact Hall atlas obstruction},\\
\texttt{/Users/raeez/calabi-yau-quantum-groups/FRONTIER.md:197--205} &
\hbox{Vol.~III Hall--Drinfeld/BKM comparison remains open}.
\end{array}
\]

\subsection{What determines primitive brackets}

A factorization or Hall object determines primitive Lie brackets only when
the chain-level multiplication and primitive projection are present.  For a
compact Hall realization this means extension correspondences
\[
\mathfrak M_\gamma^\sigma\times\mathfrak M_\eta^\sigma
\xleftarrow{\ p_{\gamma,\eta}\ }
\mathfrak E_{\gamma,\eta}^\sigma
\xrightarrow{\ q_{\gamma,\eta}\ }
\mathfrak M_{\gamma+\eta}^\sigma ,
\]
vanishing-cycle coefficients, Thom--Sebastiani orientation transport, and
proper pull-push maps
\[
m_{\gamma,\eta}=q_{\gamma,\eta!}p_{\gamma,\eta}^{*}.
\]
After a coproduct or factorization coalgebra has supplied a primitive
projection, the primitive bracket is the supercommutator of this product:
\[
[u_{\gamma,a},u_{\eta,b}]_{\Hall}
=
\sum_c C^c_{(\gamma,a),(\eta,b)}u_{\gamma+\eta,c}.
\]
The constants \(C^c_{(\gamma,a),(\eta,b)}\) depend on the extension stack,
orientation cocycle, protected integration, and chosen primitive bases.  The
Euler number
\[
\sum_a(-1)^{|a|}=\sdim \Prim_{\mathrm{prot},\gamma}
\]
does not contain any of this bilinear data.

In a curve-level chiral factorization object the same statement is phrased
as collision data.  The bracket is extracted from the singular/collision
part of the factorization product on disjoint opens, after passing to
protected primitives.  Beilinson--Drinfeld's chiral enveloping construction
constructs a chiral algebra from a supplied Lie-* algebra; it does not
construct the compact \(K3\times E\) Lie-* algebra from a scalar
determinant.

Thus the determinant controls only a Grothendieck class:
\[
[H^0\Prim_{\mathrm{prot}}(\mathcal F_X)_\gamma]\in K_0(\mathrm{sVect}),
\qquad
\sdim=f(nm,l).
\]
It does not determine products, coproducts, commutators, Jacobi constants,
or the radical quotient of a Hall--Drinfeld double.

\subsection{Borcherds--Gritsenko--Nikulin datum}

The denominator algebra used in the manuscript is fixed by the
Gritsenko--Nikulin automorphic correction \cite{GN}.  The real roots are
\[
\delta_1=2f_2-f_3,\qquad
\delta_2=2f_{-2}-f_3,\qquad
\delta_3=f_3,
\]
with Gram matrix
\[
A=((\delta_i,\delta_j))=
\begin{pmatrix}
2&-2&-2\\
-2&2&-2\\
-2&-2&2
\end{pmatrix}.
\]
The Weyl vector is
\[
\rho=\frac12(\delta_1+\delta_2+\delta_3)
=f_2-\frac12f_3+f_{-2},
\qquad
(\rho,\delta_i)=-1 .
\]
The charge-to-root map is
\[
\alpha(n,l,m)=2nf_2-lf_3+2mf_{-2}
=n\delta_1+m\delta_2+(n+m-l)\delta_3,
\]
and
\[
(\alpha(n,l,m),\alpha(n,l,m))=-2(4nm-l^2).
\]

Gritsenko--Nikulin's product theorem gives
\[
\frac1{64}\Delta_5(Z)
=
q^{1/2}r^{1/2}s^{1/2}
\prod_{(n,l,m)>0}
(1-q^nr^ls^m)^{f(nm,l)}
\]
and the Weyl--Kac--Borcherds denominator normalization used in the paper is
\[
\operatorname{den}(\autcor)
=
\frac1{64}\Delta_5(2Z)
=
e^{-2\pi i(\rho,z)}
\prod_{\alpha\in\mathcal R_+}
\left(1-e^{-2\pi i(\alpha,z)}\right)^{\smult(\alpha)}.
\]
The visible active multiplicities satisfy
\[
\smult(\alpha(n,l,m))=f(nm,l).
\]

The presentation contains more than these exponents.  One chooses an index
set \(I\) over the real and imaginary simple roots, with parity map
\(p:I\to\Z/2\Z\).  The real roots are even.  The imaginary simple roots are
read from the correction character
\[
1-\sum_{a\in\La^{2,1}_{II}\cap\R_{>0}\mathcal P_{II}}
m(a)e^{-2\pi i(a,z)}
\]
together with the isotropic ray coefficients
\[
1-\sum_{t\ge1}m(ta_0)q^t
=\prod_{t\ge1}(1-q^t)^{\tau(ta_0)},\qquad
\tau(ta_0)=9.
\]
The parity split is the Gritsenko--Nikulin BKM datum, not a consequence of
the scalar product alone:
\[
\smult(\alpha)
=\dim\mathfrak g_{\alpha,\overline0}
-\dim\mathfrak g_{\alpha,\overline1}.
\]

The generators are \(h_i,e_i,f_i\) for \(i\in I\), with
\[
[h_i,e_j]=(\alpha_i,\alpha_j)e_j,\qquad
[h_i,f_j]=-(\alpha_i,\alpha_j)f_j,\qquad
[e_i,f_j]=\delta_{ij}h_i,
\]
the real Serre relations
\[
(\operatorname{ad}e_i)^{1-2(\alpha_i,\alpha_j)/(\alpha_i,\alpha_i)}e_j=0
\quad(\alpha_i\ \mathrm{real},\ i\ne j),
\]
and the orthogonality relations
\[
(\alpha_i,\alpha_j)=0
\quad\Longrightarrow\quad
[e_i,e_j]=[f_i,f_j]=0.
\]
These are the Borcherds/Kac defining relations used by
Gritsenko--Nikulin, Appendix~\S6; Borcherds' 1988 definition gives the
same generators-and-relations universal object.

\subsection{Uniqueness: what is true}

There is no uniqueness theorem saying that a denominator identity, signed
root multiplicities, and real-root Serre relations force a compact Hall
primitive algebra to be \(\autcor\).

The available uniqueness theorem is the presentation theorem.  Fix the full
Borcherds-Cartan superdatum
\[
\mathscr B_{\Delta_5}
=
\left(
\La^{2,1}_{II}\otimes\C,\ I,\ p,\ i\mapsto\alpha_i,
(\alpha_i,\alpha_j)
\right),
\]
including the multiplicity fibres over imaginary simple roots and their
parities.  The generalized Borcherds--Kac--Moody superalgebra
\(G(\mathscr B_{\Delta_5})\) is the universal Lie superalgebra generated
by \(H\) and \(e_i,f_i\) with the relations above, modulo the standard
radical/maximal ideal required in the Kac formulation.  It is unique by
this universal property.

Consequently, a compact primitive Lie superalgebra \(P\) is identified with
\(\autcor\) if the following data are supplied:
\begin{enumerate}
\item a Cartan identification
\[
P_0\simeq\La^{2,1}_{II}\otimes\C;
\]
\item primitive classes \(e_i^P,f_i^P\) for every real and imaginary simple
root index \(i\in I\), preserving parity and multiplicity fibres;
\item the Chevalley, real Serre, and orthogonality relations above inside
the Hall/factorization bracket;
\item generation of \(P\) by these primitive classes and the Cartan;
\item the same radical quotient, or equivalently a proof that the
presentation map
\[
G(\mathscr B_{\Delta_5})\longrightarrow P
\]
has zero kernel after the completed Hall--Drinfeld pairing quotient;
\item equality of full root-space parity dimensions, not only signed
Euler characteristics.
\end{enumerate}
Under these hypotheses the induced map
\[
\autcor=G(\mathscr B_{\Delta_5})
\xrightarrow{\ \sim\ }
P
\]
is an isomorphism of charge-graded Lie superalgebras.  This is a
presentation-level recognition theorem.  It is not a consequence of the
determinant.

Even full root dimensions would not by themselves determine the brackets.
A graded Hilbert series is not a multiplication table.  The first concrete
place where the missing constants appear is the degree
\[
\delta_{123}:=\delta_1+\delta_2+\delta_3=\alpha(1,1,1).
\]
For \(a_{ij}=\delta_i+\delta_j\), the denominator algebra has one
real-real commutator direction \(E_{ij}=[e_i,e_j]\) and nine isotropic
simple directions \(u_{ij,r}\), so
\[
\smult(a_{ij})=10=1+9.
\]
The brackets
\[
[e_k,u_{ij,r}]\in\mathfrak g_{\delta_{123}}
\qquad(\{i,j,k\}=\{1,2,3\})
\]
are not numbers until a basis and parity split of
\(\mathfrak g_{\delta_{123}}\) have been supplied.  The determinant gives
only
\[
\smult(\delta_{123})=f(1,1)=-64.
\]
It gives no basis, no parity split, and no expression of
\([e_k,u_{ij,r}]\) in that basis.

\subsection{Attack on determinant-to-algebra}

The determinant-to-algebra implication fails for elementary reasons.

First, let
\[
B=\bigoplus_{\gamma\in\Gamma_{\mathrm{eff}}}B_\gamma
\]
be any charge-graded finite-dimensional super vector space with
\[
\sdim B_{(n,l,m)}=f(nm,l).
\]
Then
\[
64q^{1/2}r^{1/2}s^{1/2}
\prod_{\gamma\in\Gamma_{\mathrm{eff}}}(1-x^\gamma)^{\sdim B_\gamma}
=\Delta_5(Z).
\]
Putting the zero bracket on \(B\) gives the same determinant.  It is not
\(\autcor\), since in the denominator algebra
\[
[e_i,e_j]=E_{ij}\ne0\qquad(i\ne j).
\]

Second, even after fixing \(\autcor\), let \(M\) be any nonzero
root-graded \(\autcor\)-module with finite-dimensional weight spaces.
The abelian ideal
\[
M\oplus \Pi M
\]
has zero signed character.  The semidirect product
\[
\autcor\ltimes (M\oplus\Pi M)
\]
has the same signed denominator as \(\autcor\), but is not isomorphic to
\(\autcor\).  Thus signed denominator data do not exclude invisible
zero-superdimension root spaces or abelian ideals.
This is exactly the warning in the manuscript at
\texttt{main.tex:2274--2277} and \texttt{main.tex:2291--2303}.

Third, the OP/DT scalar square is still farther from the bracket.  It
computes a reduced scalar partition function
\[
Z^X_{\mathrm{OP}}=-4096\,\Delta_5^{-2}.
\]
It does not construct an extension-closed heart, an oriented critical Hall
atlas, an \(E\)-equivariant/reduced Hall convolution, or a primitive
projection.  Therefore it supplies no Hall constants.

\subsection{Strongest conditional theorem}

\begin{theorem}[Conditional primitive comparison]
\label{thm:conditional-primitive-comparison-gdelta5}
Let \(X=K3\times E\).  Suppose the following objects have been constructed.
\begin{enumerate}
\item A charge-completed holomorphic \(E_3\)-factorization algebra
\(\mathcal A_X^{E_3}\) on \(X\), finite-first after sectorial
Harder--Narasimhan truncation.
\item A chain-level \(K3\)-to-\(E\) specialization
\[
\mathcal F_X\simeq
\mathrm{Sp}^{\mathrm{ch}}_{K3,E}(\mathcal A_X^{E_3})
\in\Fact(\operatorname{Ran}(E)).
\]
\item An oriented compact critical Hall factorization object
\(\Hall_X\) with extension correspondences, vanishing cycles,
Thom--Sebastiani orientation transport, support property, HN descent, and
an \(E\)-equivariant or reduced convolution retaining relative
\(E\)-position.
\item A protected primitive functor and supercommutator
\[
[-,-]_{\mathrm{prot}}:
H^0\Prim_{\mathrm{prot}}(\mathcal F_X)_\gamma
\otimes
H^0\Prim_{\mathrm{prot}}(\mathcal F_X)_\eta
\to
H^0\Prim_{\mathrm{prot}}(\mathcal F_X)_{\gamma+\eta}.
\]
\item A protected integration and parity splitting identifying the full
Gritsenko--Nikulin root-space parity dimensions, and in particular
\[
\sdim H^0\Prim_{\mathrm{prot}}(\mathcal F_X)_{(n,l,m)}
=f(nm,l).
\]
\item An orientation character satisfying
\[
\epsilon_o|_{W^{(2)}(\La^{2,1}_{II})}
=
\nu_{\Delta_5}|_{W^{(2)}(\La^{2,1}_{II})}.
\]
\item Primitive representatives for every real and imaginary simple root
of \(\mathscr B_{\Delta_5}\), preserving multiplicities, parities, and
the degree map \(\alpha(n,l,m)=2nf_2-lf_3+2mf_{-2}\).
\item A proof that the protected primitive bracket satisfies the full
Borcherds/Kac defining relations for \(\mathscr B_{\Delta_5}\): Cartan,
Chevalley, real Serre, imaginary orthogonality, parity signs, and the
same completed Hopf pairing and radical quotient.
\item Generation of
\(H^0\Prim_{\mathrm{prot}}(\mathcal F_X)\) by these simple primitive
classes and the Cartan.
\end{enumerate}
Then
\[
H^0\Prim_{\mathrm{prot}}(\mathcal F_X)
\cong
\autcor
\]
as charge-graded Lie superalgebras.  Its Weyl--Kac--Borcherds
denominator is \(64^{-1}\Delta_5(2Z)\).
\end{theorem}

\begin{proof}
The listed simple primitive classes and relations give a homomorphism from
the universal Borcherds--Kac--Moody superalgebra
\(G(\mathscr B_{\Delta_5})=\autcor\) to
\(H^0\Prim_{\mathrm{prot}}(\mathcal F_X)\).  Generation gives
surjectivity.  Equality of full root-space parity dimensions together with
the same completed radical quotient gives injectivity on every root space
and on the Cartan.  Hence the homomorphism is an isomorphism.  The
denominator statement is then Theorem~\ref{thm:denominator-algebra-identity}.
\end{proof}

\subsection{Claim-status recommendation}

The manuscript should keep
\[
H^0\Prim_{\mathrm{prot}}(\mathcal F_X)
\stackrel{?}{\cong}
\autcor
\]
as the compact \(E_3\)/Hall/chiral realization problem.  It may state
Theorem~\ref{thm:conditional-primitive-comparison-gdelta5} as a conditional
recognition criterion.  It must not say that the determinant, the scalar
square, or the denominator identity proves the primitive comparison.

The precise missing geometric hypotheses are:
\[
\begin{gathered}
\mathcal A_X^{E_3},\quad
\mathrm{Sp}^{\mathrm{ch}}_{K3,E},\quad
\mathcal F_X,\quad
\Hall_X,\quad
\mathfrak M_\gamma^\sigma,\quad
\mathfrak E_{\gamma,\eta}^\sigma,\quad
o,\quad
\epsilon_o=\nu_{\Delta_5},\\
\Prim_{\mathrm{prot}},\quad
\chi^{\mathrm{prot}},\quad
\hbox{full parity split},\quad
\hbox{simple primitive representatives},\quad
\hbox{Borcherds relations in the Hall bracket},\quad
\hbox{completed Hopf pairing/radical quotient}.
\end{gathered}
\]

\subsection{Source anchors}

Primary source anchors:
\begin{enumerate}
\item Gritsenko--Nikulin \cite{GN}, Introduction and Theorem~4.1:
\(\Delta_5\) is the correcting Siegel modular form for the rank-three
matrix \(G_1=A\), and
\[
64^{-1}\Delta_5(Z)
=e^{\pi i(z_1+z_2+z_3)}
\prod_{(n,l,m)>0}
\left(1-e^{2\pi i(nz_1+lz_2+mz_3)}\right)^{f(nm,l)}.
\]
\item Gritsenko--Nikulin \cite{GN}, Appendix~\S6: generalized
Kac--Moody superalgebras without odd real simple roots have the
Chevalley, real Serre, and orthogonality relations used in
\texttt{main.tex:2149--2174}.
\item Borcherds, \emph{Generalized Kac--Moody algebras}, Journal of
Algebra 115 (1988), \S1: the algebra associated to a symmetrized
Cartan datum is generated by \(H,e_i,f_i\) with relations (1)--(5).
\item Borcherds \cite{B}, Theorem~13.3: the singular theta lift gives
automorphic forms with infinite product expansions; at singular weight
these products supply denominator formulae for generalized Kac--Moody
superalgebras.
\item Kontsevich--Soibelman \cite{KSCoHA} and Davison--Meinhardt
\cite{DavisonMeinhardtCoHA}: Hall products require critical stacks,
vanishing cycles, orientations, and extension correspondences.  These
are the sources of Hall structure constants, not the scalar determinant.
\item Beilinson--Drinfeld \cite{BD} and Costello--Gwilliam
\cite{CostelloGwilliam1,CostelloGwilliam2}: factorization/chiral
products are chain-level local operations.  A global trace or determinant
is their decategorified shadow.
\end{enumerate}

\subsection{Files changed, checks, open problems}

Files changed:
\[
\texttt{agent\_material/12\_primitive\_comparison\_attack\_heal.tex}.
\]

Checks run:
\[
\begin{array}{l}
\texttt{rg -n "denominator|Hall|factorization|determinant|Prim|primitive" main.tex},\\
\texttt{nl -ba main.tex | sed -n '230,310p'},\\
\texttt{nl -ba main.tex | sed -n '1560,2100p'},\\
\texttt{nl -ba main.tex | sed -n '2110,2385p'},\\
\texttt{nl -ba agent\_material/10\_bps\_bracket\_constants\_attack\_heal.tex},\\
\texttt{nl -ba agent\_material/11\_factorization\_koszul\_minimal\_data.tex},\\
\texttt{nl -ba agent\_material/08\_compact\_hall\_factorization\_object.tex},\\
\texttt{nl -ba /Users/raeez/calabi-yau-quantum-groups/FRONTIER.md | sed -n '1,260p'},\\
\texttt{rg -n <forbidden-register-scan> agent\_material/12\_primitive\_comparison\_attack\_heal.tex},\\
\texttt{rg -n "\textbackslash\textbackslash(begin|end|cite|ref)" agent\_material/12\_primitive\_comparison\_attack\_heal.tex},\\
\texttt{git diff --no-index -- /dev/null agent\_material/12\_primitive\_comparison\_attack\_heal.tex},\\
\texttt{git status --short -- agent\_material/12\_primitive\_comparison\_attack\_heal.tex}.
\end{array}
\]
No manuscript build was run.  This file is standalone agent material and
is not included from \texttt{main.tex}.

Open problems:
\begin{enumerate}
\item Construct the compact \(E_3\)-factorization algebra
\(\mathcal A_X^{E_3}\) and its chain-level specialization to \(E\).
\item Construct the compact oriented Hall atlas with HN finite-first
completion and reduced/equivariant \(E\)-descent.
\item Prove the orientation character comparison
\(\epsilon_o=\nu_{\Delta_5}\).
\item Construct protected primitives and prove their full parity
dimensions, not only their signed Euler characteristics.
\item Identify primitive representatives of all real and imaginary
simple roots, including the \(9\) isotropic simple directions on each
primitive isotropic ray.
\item Prove the Chevalley, real Serre, orthogonality, and imaginary-root
relations inside the Hall/factorization supercommutator.
\item Construct the continuous Hopf pairing, negative half, Cartan
completion, and radical quotient giving the Hall--Drinfeld double.
\end{enumerate}
